% 第 15 章 Creating Anchors
\bichapter{创建 Anchor}{Creating Anchors}
\label{chap:gui-anchor}

Anchor 把一个设计对象固定到另一个对象上。被固定的对象是从属对象（dependent object），它所依附的对象是锚点对象（anchor object）。移动锚点对象时，固定在它上面的从属对象一起移动。移动从属对象时，它会弹回锚点位置。

在 IC Compiler II GUI 中选择 \textbf{Create > Anchor}，即可打开创建工具。该工具提供：
\begin{itemize}
  \item Anchor type：指定固定方式，可以只固定 X、只固定 Y，或同时固定 X 与 Y。
  \item Dependent：指定从属对象。
  \item Anchor：指定锚点对象。
  \item Position Reference：指定从属对象与锚点对象的位置参考。
  \item Offset：指定从属对象相对锚点对象的 X、Y 偏移，默认是 \cmd{\{0,0\}}。
\end{itemize}

也可以在层次内部创建 anchor。在 Layout settings 面板中启用 deep selection 并提高 view level，然后即可在层次内部点选从属对象和锚点对象。

\begin{noteBox}
从属对象的层次级别必须等于或低于锚点对象的层次级别。
\end{noteBox}

表~\ref{tab:gui-anchor} 是 CreateAnchor 工具中交互编辑的约定。

\begin{table}[htbp]
\centering
\caption{CreateAnchor 工具的交互约定}
\label{tab:gui-anchor}
\begin{tabularx}{\textwidth}{@{}l X@{}}
\toprule
事件 & 动作 \\
\midrule
单击或拖拽 & 设置矩形一角，或创建 Anchor \\
双击或按 Enter & 在当前位置创建 Anchor \\
中键单击 & 激活笔势菜单 \\
右键单击 & 显示上下文菜单 \\
F1 & 在重叠对象中向前循环 \\
Shift+F1 & 在重叠对象中向后循环 \\
Escape & 若有未完成的更改则取消；否则退出工具 \\
\bottomrule
\end{tabularx}
\end{table}
